%----------------------------------------------------------------------------------------
% HET BOEK: het documentskelet dat alle cursussen delen. Een cursus levert alleen
%   \cursustitel         de titel
%   \cursushoofdstukken  de lijst met \hoofdstuk{nummer}{bestand}
% en zijn eigen notatie.tex, voorkant-gegevens.tex, hoofdstukken/ en oefeningen/.
% Pas dit bestand alleen hier aan: het geldt voor alle cursussen.
%----------------------------------------------------------------------------------------
\documentclass[11pt,fleqn,oneside,openany]{book}

\newcommand{\cursusondertitel}{Handboek}
\newcommand{\cursusauteur}{Axiomath}

% De voorkant: gegevens en motief van deze cursus (Ontwerp/voorkanten/maak-voorkanten.py).
\input{voorkant-gegevens}

\input{\hoofdmap Opmaak/opmaak}
\input{notatie}

\ifleerkracht\renewcommand{\cursusondertitel}{Leerkrachtenversie}\fi
\ifzelfstudie\renewcommand{\cursusondertitel}{Zelfstudie voor thuis}\fi

\providecommand{\uitvoer}{cursus}
\newif\ifcursus
\ifdefstring{\uitvoer}{cursus}{\cursustrue}{\cursusfalse}

% Zie maak-pdf.py: gen/ bevat de hoofdstukken met de oefeningen uit de bank en de bewaarde
% grafieken. Draai dat script dus eerst. \enkelhoofdstuk bouwt maar één hoofdstuk.
\providecommand{\hoofdstukmap}{gen}
\newcommand{\hoofdstuk}[2]{%
  \ifcsname enkelhoofdstuk\endcsname
    \ifnum#1=\enkelhoofdstuk\relax\setcounter{chapter}{\numexpr#1-1\relax}\input{\hoofdstukmap/#2}\fi
  \else
    \input{\hoofdstukmap/#2}%
  \fi}

\begin{document}

\ifcursus
\Opensolutionfile{oplossingen}
\Opensolutionfile{uitgewerkteoplossingen}
\fi

\ifcursus\else
\renewcommand{\cursusondertitel}{\ifdefstring{\uitvoer}{numeriek}{Numerieke oplossingen}{Uitgewerkte oplossingen}}
\fi
\maakomslag

\ifcursus

\inhoudstafel
\input{\hoofdmap Opmaak/zo-werk-je}
\clearpage % voor de teller van het hoofdstuk verandert: deze pagina krijgt geen tab

\cursushoofdstukken

\Closesolutionfile{oplossingen}
\Closesolutionfile{uitgewerkteoplossingen}

% De leerkrachten- en zelfstudieversie krijgen de antwoorden meteen achteraan
\ifuitwerking
\clearpage\global\opdrachtdeelfalse
\section*{Antwoorden}
\phantomsection\addcontentsline{toc}{chapter}{Antwoorden}
\begin{multicols}{2}\footnotesize
\input{oplossingen}
\end{multicols}
\fi

\else

\ifdefstring{\uitvoer}{numeriek}{%
\chapter*{Numerieke oplossingen}
\begin{multicols}{2}
\input{oplossingen}
\end{multicols}
}{%
\chapter*{Uitgewerkte oplossingen}
\input{uitgewerkteoplossingen}
}

\fi

\end{document}
